\subsection{基本ソフトウェア}

基本ソフトウェア、すなわちオペレーティングシステム（OS）は、コンピュータのハードウェアを管理し、アプリケーションに実行基盤を提供するソフトウェアである。OS の理解は、性能、並行性、入出力、ファイル、セキュリティ、ネットワークの設計判断に直結する。

OS には、バッチ処理、マルチプログラミング、時分割、単一利用者、複数利用者、単一タスク、マルチタスク、マルチスレッド、リアルタイム、ネットワーク、分散などの種類がある。小規模な組込みシステムでは、OS を使わない構成もあり得る。

\par\smallskip\noindent\refstepcounter{table}\label{tab:ch16-16}表 \thetable: 基本ソフトウェアにおけるOSの主要管理対象・説明\par\smallskip
表 \thetable は、基本ソフトウェアにおけるOSの主要管理対象・説明を比較・確認しやすい形で整理したものである。\par\smallskip
\begin{longtable}{>{\raggedright\arraybackslash}p{0.28\textwidth}>{\raggedright\arraybackslash}p{0.66\textwidth}}
\toprule
OSの主要管理対象 & 説明 \\
\midrule
\endfirsthead
\toprule
OSの主要管理対象 & 説明 \\
\midrule
\endhead
プロセッサ管理 & プロセス、スレッド、スケジューリング、同期、デッドロックを扱う。 \\
メモリ管理 & 物理メモリ、仮想メモリ、ページング、置換、断片化を扱う。 \\
装置管理 & 入出力装置、装置ドライバ、割り込み、DMA、バッファリングを扱う。 \\
情報管理 & ファイル、ディレクトリ、アクセス制御、記憶装置、保護を扱う。 \\
ネットワーク管理 & 分散処理、同期、相互排除、時刻、通信、セキュリティを扱う。 \\
\bottomrule
\end{longtable}

\par\smallskip
\begin{center}
\centering
\IfFileExists{assets/generated_figures/ch16/fig-ch16-11.png}{\includegraphics[width=0.96\linewidth]{assets/generated_figures/ch16/fig-ch16-11.png}}{\fbox{\parbox[c][48mm][c]{0.9\linewidth}{\centering 画像生成待ち\\OSが管理する資源の全体図}}}
\par\smallskip
\refstepcounter{figure}\label{fig:ch16-11}図 \thefigure: OSが管理する資源の全体図
\end{center}
図 \thefigure は、OSが管理する資源の全体図を視覚的に整理したものである。
\par\smallskip

\subsubsection{プロセッサ管理}

プロセッサ管理では、プロセス、アドレス空間、スレッド、同期、ロック、スケジューリングを理解する必要がある。プロセスは実行中のプログラムの単位であり、スレッドはプロセス内の実行の流れである。複数の処理が同時に動くと、共有データへのアクセス順序が問題になる。

プロセス間通信には、メッセージ、パイプ、共有メモリ、セマフォなどがある。同期が不適切だと、競合、デッドロック、飢餓が発生する。CPU スケジューリングには、到着順、最短ジョブ優先、最短残り時間優先、優先度、ラウンドロビンなどがある。

\subsubsection{メモリ管理}

メモリ管理では、物理メモリ、仮想メモリ、二次記憶、記憶階層、リンク、割り当てを理解する。断片化には、未使用領域が分散する外部断片化と、割り当て単位の内部に無駄が残る内部断片化がある。ページング、セグメンテーション、デマンドページング、ページ置換、スラッシング、スワッピングは重要な概念である。

ページ置換戦略には、FIFO、NRU、LRU、MRU、LFU、MFU、セカンドチャンス、エージングなどがある。アプリケーションのメモリ使用特性が悪いと、OS の工夫だけでは性能を回復できないことがある。

\subsubsection{装置管理}

装置管理では、メモリマップ入出力と入出力マップ入出力、ブロック装置と文字装置、ポーリング、割り込み駆動、DMA、ブロッキング入出力と非ブロッキング入出力を理解する。装置ドライバは、ハードウェアとアプリケーションの間に立つソフトウェアである。

\subsubsection{情報管理}

情報管理では、ファイルシステム、記憶管理、ファイル属性、ディレクトリ構造、アクセス制御、保護、セキュリティを扱う。アクセス制御リスト、アクセス行列、権限、認証、ハードリンク、ソフトリンクなどは、ソフトウェアの安全性と運用性に関係する。

\subsubsection{ネットワーク管理}

ネットワーク管理では、分散オブジェクト、分散ファイルシステム、同期、時刻、分散相互排除、選挙アルゴリズム、マルチキャスト通信などを扱う。分散環境では、単一の正しい時刻を前提にしにくく、論理時計やベクトル時計のような概念が必要になる。

\begin{pointbox}{実務での注意}
基本ソフトウェアの制約は、アプリケーションの設計制約でもある。スレッド、ファイル、メモリ、ネットワーク、権限の扱いを知らずに実装すると、性能問題、データ破壊、セキュリティ事故、再現困難な障害につながる。

\end{pointbox}
